%% ----------------------------------------------------------------------------
%%  §4.8  Appendice: il metodo di bisezione
%% ----------------------------------------------------------------------------
\section{Appendice: il metodo di bisezione per il calcolo delle radici}
\label{sec:bisezione}

In questo paragrafo dimostriamo il teorema dell'esistenza degli zeri (\cref{teo:zeri}).
Utilizziamo nella dimostrazione il \emph{metodo di bisezione}, già incontrato nel
\cref{cap:numeri} per approssimare $\sqrt2$ e nel \cref{cap:successioni} per dimostrare
il teorema di Bolzano--Weierstrass: si tratta di un procedimento costruttivo che, oltre a
dimostrare l'esistenza di una soluzione di un'equazione data, fornisce anche un metodo
per calcolarla.

Prendiamo in considerazione \emph{equazioni} del tipo
\begin{equation}\label{eq:equazione}
  f(x)=0,
\end{equation}
con $f$ funzione definita in un intervallo $[a,b]$. Risolvere l'equazione significa
determinare tutti i numeri reali $x_{0}\in[a,b]$ per cui $f(x_{0})=0$; tali numeri si
dicono \emph{soluzioni} dell'equazione \eqref{eq:equazione}, o anche \emph{zeri} della
funzione $f$. Se $f$ è un polinomio, si dice che \eqref{eq:equazione} è un'\emph{equazione
algebrica}; se $f$ è una funzione trascendente (ad esempio composta tramite le funzioni
$e^{x}$, $\log x$, $\sen x$, $\cos x$), la \eqref{eq:equazione} prende il nome di
\emph{equazione trascendente}. Una soluzione di un'equazione algebrica si dice anche
\emph{radice} dell'equazione; per estensione, si usa il termine di radici anche per le
soluzioni di equazioni trascendenti.

\begin{strategia}[Piano d'attacco per il \cref{teo:zeri}]
Ricordiamo le ipotesi: $f$ è continua in $[a,b]$ e $f(a)<0<f(b)$. L'idea è di
\emph{intrappolare} uno zero in intervalli sempre più piccoli, mantenendo l'invariante
«$f$ è negativa nell'estremo sinistro e positiva nell'estremo destro». Si dimezza
l'intervallo e si tiene la metà agli estremi della quale $f$ continua a cambiare segno.
Gli estremi sinistri formano una successione crescente e limitata, che converge per il
teorema sulle successioni monotòne (qui entra l'assioma di completezza); gli estremi
destri convergono allo stesso limite $x_{0}$, perché le lunghezze si dimezzano. Infine
la continuità di $f$ permette di passare al limite nelle disuguaglianze $f(a_{n})<0$ e
$f(b_{n})>0$, che diventano $f(x_{0})\le0$ e $f(x_{0})\ge0$.
\end{strategia}

\begin{dimostrazione}[Dimostrazione del \cref{teo:zeri}]
\emph{Costruzione.} Consideriamo il numero $c$, punto di mezzo dell'intervallo $[a,b]$,
cioè $c=\frac{a+b}2$. Se $f(c)=0$ abbiamo trovato una radice. Altrimenti consideriamo i
due casi $f(c)>0$, $f(c)<0$. Se $f(c)>0$, la funzione $f$ assume valori di segno discorde
agli estremi dell'intervallo $[a,c]$, mentre se $f(c)<0$ è $[c,b]$ l'intervallo dove $f$
cambia segno. Indichiamo con $[a_{1},b_{1}]$ l'intervallo da considerare, cioè definiamo
\begin{equation}\label{eq:bisezione-primo}
  \begin{cases}
    \text{se } f(c)>0 \ \implica\ a_{1}=a,\ b_{1}=c;\\
    \text{se } f(c)<0 \ \implica\ a_{1}=c,\ b_{1}=b.
  \end{cases}
\end{equation}
Abbiamo così trovato un intervallo $[a_{1},b_{1}]$, di ampiezza metà del precedente
$[a,b]$, per cui risulta $f(a_{1})<0$, $f(b_{1})>0$. Definiamo $c_{1}=\frac{a_{1}+b_{1}}2$ e
ripetiamo il ragionamento. Posto $a_{0}=a$, $b_{0}=b$, otteniamo tre successioni $a_{n}$,
$b_{n}$, $c_{n}$ definite per ricorrenza da $c_{n}=\frac{a_{n}+b_{n}}2$ e
\begin{equation}\label{eq:bisezione-ricorrenza}
  \begin{cases}
    \text{se } f(c_{n})>0 \ \implica\ a_{n+1}=a_{n},\ b_{n+1}=c_{n};\\
    \text{se } f(c_{n})<0 \ \implica\ a_{n+1}=c_{n},\ b_{n+1}=b_{n}.
  \end{cases}
\end{equation}
Se per qualche $n$ risulta $f(c_{n})=0$, ci si ferma perché si è trovata una radice, e
il teorema è dimostrato. Altrimenti le successioni sono definite per ogni $n$ e, per
costruzione, risulta
\begin{equation}\label{eq:bisezione-segni}
  f(a_{n})<0, \qquad f(b_{n})>0 \qquad\perogni n\ge0 .
\end{equation}

\emph{Lunghezze.} È semplice scrivere la relazione che lega $a_{n}$ con $b_{n}$: infatti,
a ogni iterazione, la lunghezza dell'intervallo $[a_{n},b_{n}]$ si dimezza. Quindi
$b_{1}-a_{1}=\frac{b-a}2$, $b_{2}-a_{2}=\frac{b-a}{2^{2}}$ e, dopo $n$ passi (per
induzione),
\begin{equation}\label{eq:bisezione-lunghezza}
  b_{n}-a_{n}=\frac{b-a}{2^{n}} \qquad\perogni n\ge0 .
\end{equation}

\emph{Convergenza.} Per costruzione, la successione $a_{n}$ è crescente
($a_{0}\le a_{1}\le a_{2}\le\cdots$), la successione $b_{n}$ è decrescente, e
$a\le a_{n}<b_{n}\le b$ per ogni $n$: in particolare $a_{n}$ è limitata. Per il teorema
sulle successioni monotòne (\cref{teo:monotone}) $a_{n}$ ammette limite finito; sia
$x_{0}$ tale limite. Anche la successione $b_{n}$, espressa mediante la
\eqref{eq:bisezione-lunghezza} da
\begin{equation}\label{eq:bisezione-bn}
  b_{n}=a_{n}+\frac{b-a}{2^{n}},
\end{equation}
converge a $x_{0}$ per $n\to+\infty$, perché $\frac1{2^{n}}\to0$ (\cref{teo:lim-geometrica}).
Inoltre $x_{0}\in[a,b]$, per il \cref{cor:confronto-limiti} applicato alle
disuguaglianze $a\le a_{n}\le b$.

\emph{Passaggio al limite.} Quindi, ricordando la \eqref{eq:bisezione-segni}, dalla
continuità di $f$ nel punto $x_{0}$ (\cref{prop:caratt-continuita}(b)) e dal
\cref{cor:limite-non-negativo} si ottiene
\begin{equation}\label{eq:bisezione-limite}
  f(x_{0})=\lim_{n\to+\infty}f(a_{n})\le0, \qquad f(x_{0})=\lim_{n\to+\infty}f(b_{n})\ge0 .
\end{equation}
Perciò $f(x_{0})=0$. Infine $x_{0}\neq a$ e $x_{0}\neq b$, perché $f(a)<0$ e $f(b)>0$;
dunque $x_{0}\in(a,b)$, e il teorema dell'esistenza degli zeri è provato.
\end{dimostrazione}

\begin{strategia}[Un'altra via]
Invece di costruire lo zero, lo si individua con l'estremo superiore: si considera
l'insieme $E$ dei punti di $[a,b]$ in cui $f$ è negativa, e il suo estremo superiore $y$.
Se fosse $f(y)<0$, oppure $f(y)>0$, la permanenza del segno darebbe un intorno di $y$ in
cui $f$ conserva il segno, e in entrambi i casi $y$ non potrebbe essere il minimo dei
maggioranti di $E$. L'assioma di completezza interviene direttamente, nella forma
dell'esistenza dell'estremo superiore; la continuità, soltanto attraverso la permanenza
del segno.
\end{strategia}

\begin{dimostrazione}[Seconda dimostrazione]
Sia
\begin{equation}\label{eq:zeri-insieme-E}
  E=\{x\in[a,b] : f(x)<0\} .
\end{equation}
L'insieme $E$ non è vuoto, perché $a\in E$, ed è limitato superiormente da $b$; per il
\cref{teo:esistenza-sup} esiste $y=\sup E$, e $a\le y\le b$. Dimostriamo che $f(y)=0$,
escludendo gli altri due casi.

\emph{Non può essere $f(y)<0$.} In tal caso $y\neq b$, perché $f(b)>0$; dunque $y<b$. La
funzione $-f$ è continua in $y$ (\cref{teo:operazioni-continue}) e $-f(y)>0$: per la
permanenza del segno (\cref{teo:permanenza-segno-continue}) esiste $\delta>0$ tale che
$f(x)<0$ per ogni $x\in[a,b]$ con $\abs{x-y}<\delta$. Il punto
$x_{1}=\min\bigl\{y+\frac\delta2,b\bigr\}$ appartiene a $[a,b]$ e verifica
$y<x_{1}<y+\delta$; quindi $f(x_{1})<0$, cioè $x_{1}\in E$. Ma $x_{1}>y$, contro il fatto
che $y$ è un maggiorante di $E$.

\emph{Non può essere $f(y)>0$.} Per la permanenza del segno esiste $\delta>0$ tale che
$f(x)>0$ per ogni $x\in[a,b]$ con $\abs{x-y}<\delta$. Per la caratterizzazione
dell'estremo superiore (\cref{prop:caratt-sup}), con $\epsilon=\delta$, esiste $x_{2}\in E$
tale che $y-\delta<x_{2}\le y$. Allora $\abs{x_{2}-y}<\delta$, e quindi $f(x_{2})>0$;
d'altra parte $f(x_{2})<0$, perché $x_{2}\in E$: è una contraddizione.

Resta $f(y)=0$. Infine $y\neq a$ e $y\neq b$, perché $f(a)<0<f(b)$; quindi $y\in(a,b)$.
\end{dimostrazione}

\begin{osservazione}[Due dimostrazioni a confronto]\label{oss:zeri-due-vie}
Le due dimostrazioni poggiano entrambe sull'assioma di completezza, ma lo usano in modo
diverso. La bisezione passa attraverso il teorema sulle successioni monotòne ed è
\emph{costruttiva}: produce le approssimazioni $c_{n}$ e ne controlla l'errore. La
seconda usa l'esistenza dell'estremo superiore e non indica alcun modo per calcolare
$\sup E$: garantisce lo zero senza dire come avvicinarlo. In compenso ne individua uno
ben preciso: $y=\sup E$ è uno zero a destra del quale $f$ non assume più valori
negativi, perché i punti di $(y,b]$ non stanno in $E$. Infine, la seconda dimostrazione
usa la continuità solo «localmente», attraverso la permanenza del segno nel punto $y$:
la stessa struttura ritorna in molte dimostrazioni di esistenza dell'Analisi.
\end{osservazione}

\begin{notastorica}[Bolzano, Cauchy e il teorema degli zeri]
Il saggio di Bolzano del 1817, ricordato nel \cref{sec:sup} per il principio
dell'estremo superiore, ha per titolo l'enunciato stesso del nostro teorema:
\emph{Dimostrazione puramente analitica del teorema che fra due valori che danno
risultati di segno opposto giace almeno una radice reale dell'equazione}. Bolzano vi
definì la continuità in modo sostanzialmente identico al nostro, e ricavò il teorema
dall'esistenza dell'estremo superiore: è l'idea della seconda dimostrazione. Quattro
anni dopo Cauchy, nel \emph{Cours d'analyse}, si accontentò nel testo di un argomento geometrico
(il grafico deve incontrare la retta), ma in una nota dedicata alla risoluzione numerica
delle equazioni diede una dimostrazione diretta: divideva l'intervallo in $m$ parti
uguali, ne sceglieva una ai cui estremi la funzione ha segni opposti e ripeteva
l'operazione, ottenendo due successioni monotòne di estremi con lo stesso limite. È la
nostra prima dimostrazione, che corrisponde a $m=2$; con $m=10$ si guadagna una cifra
decimale a ogni passo. Né l'uno né l'altro, però, poteva giustificare fino in fondo il
passo decisivo, l'esistenza dell'estremo superiore o del limite di una successione
monotòna e limitata: mancava una teoria dei numeri reali, che arrivò mezzo secolo più
tardi (si veda la nota storica del \cref{sec:bw-cauchy}).
\end{notastorica}

\begin{figure}[htb]
\centering
\includegraphics[width=0.92\textwidth]{bisezione}
\caption{Il metodo di bisezione per $f(x)=x^{3}+x-1$ su $[0,1]$: i primi intervalli
incapsulati $[a_{n},b_{n}]$ si stringono attorno alla radice $x_{0}=0{,}6823\ldots$, e a
ogni passo la loro lunghezza si dimezza (\texttt{python/fig\_bisezione.py}).}
\label{fig:bisezione}
\end{figure}

\subsection{Il calcolo numerico della radice}

Dalla dimostrazione proposta risulta chiaro come calcolare numericamente la soluzione
$x_{0}$. Infatti le tre successioni $a_{n}$, $b_{n}$, $c_{n}$ convergono a $x_{0}$ (la terza
perché è compresa tra le prime due). I termini delle tre successioni sono valori
approssimati di $x_{0}$; in particolare, i valori di $a_{n}$ sono approssimazioni per
difetto, quelli di $b_{n}$ sono approssimazioni per eccesso, cioè
\begin{equation}\label{eq:bisezione-stima}
  a_{n}\le x_{0}\le b_{n} \qquad\perogni n\ge0
\end{equation}
(la successione crescente $a_{n}$ non supera il suo limite, la decrescente $b_{n}$ non
scende sotto il suo, per il \cref{teo:monotone}). Dalle \eqref{eq:bisezione-lunghezza},
\eqref{eq:bisezione-stima} si deduce che l'errore di approssimazione che si commette
sostituendo $x_{0}$ con $a_{n}$ (oppure con $b_{n}$) non supera $\frac{b-a}{2^{n}}$. Dato
che $c_{n}$ è il punto di mezzo dell'intervallo $[a_{n},b_{n}]$, che contiene $x_{0}$,
l'errore che si commette nell'approssimare $x_{0}$ con $c_{n}$ non supera
$\frac{b-a}{2^{n+1}}$:
\begin{equation}\label{eq:bisezione-errore}
  \abs{c_{n}-x_{0}}\le\frac{b_{n}-a_{n}}2=\frac{b-a}{2^{n+1}} .
\end{equation}

Riprendiamo in considerazione le equazioni \eqref{eq:cubica}, \eqref{eq:trascendente}.
Ci proponiamo il calcolo delle rispettive radici con un errore inferiore a $10^{-3}$. In
entrambi i casi abbiamo un intervallo di ampiezza $b-a=1$; infatti nel primo caso
$[a,b]=[0,1]$, nel secondo $[a,b]=[-1,0]$. L'errore di approssimazione che si commette
sostituendo la soluzione $x_{0}$ con $c_{n}$ non supera $\frac1{2^{n+1}}$; in particolare,
per $n=9$, risulta
\begin{equation}\label{eq:bisezione-c9}
  \abs{c_{9}-x_{0}}\le\frac1{2^{10}}=\frac1{1024}<10^{-3} .
\end{equation}
Il programma seguente esegue il calcolo, e alla fine prosegue fino a $50$ dimezzamenti
per confrontare il risultato con la formula di Cardano \eqref{eq:cardano}.

\begin{python}[Il metodo di bisezione]
import math

def bisezione(f, a, b, passi):
    """Se f(a) < 0 < f(b), restituisce i punti medi c_0, ..., c_n."""
    assert f(a) < 0 < f(b)
    medi = []
    for n in range(passi + 1):
        c = (a + b) / 2
        medi.append(c)
        if f(c) == 0:          # radice esatta: ci si ferma
            break
        if f(c) > 0:
            b = c              # il cambio di segno è in [a, c]
        else:
            a = c              # il cambio di segno è in [c, b]
    return medi

f1 = lambda x: x**3 + x - 1           # f1(0) = -1 < 0 < 1 = f1(1)
f2 = lambda x: math.exp(x) + x        # f2(-1) = 1/e - 1 < 0 < f2(0)

c1 = bisezione(f1, 0, 1, 9)
c2 = bisezione(f2, -1, 0, 9)
for n in range(10):
    print(f"c_{n} = {c1[n]:.6f}     c_{n} = {c2[n]:+.6f}")

# 50 passi: errore < 2^-51; confronto con la formula di Cardano
r1 = bisezione(f1, 0, 1, 50)[-1]
s = math.sqrt(31 / 108)
cardano = (0.5 + s) ** (1/3) - (s - 0.5) ** (1/3)
print(f"x^3 + x - 1 = 0:  {r1:.15f}   (Cardano: {cardano:.15f})")
print(f"e^x + x = 0:     {bisezione(f2, -1, 0, 50)[-1]:.15f}")
\end{python}
\begin{risultato}
c_0 = 0.500000     c_0 = -0.500000
c_1 = 0.750000     c_1 = -0.750000
c_2 = 0.625000     c_2 = -0.625000
c_3 = 0.687500     c_3 = -0.562500
c_4 = 0.656250     c_4 = -0.593750
c_5 = 0.671875     c_5 = -0.578125
c_6 = 0.679688     c_6 = -0.570312
c_7 = 0.683594     c_7 = -0.566406
c_8 = 0.681641     c_8 = -0.568359
c_9 = 0.682617     c_9 = -0.567383
x^3 + x - 1 = 0:  0.682327803828020   (Cardano: 0.682327803828019)
e^x + x = 0:     -0.567143290409784
\end{risultato}

Quindi la radice dell'equazione $x^{3}+x-1=0$ soddisfa $\abs{x_{0}-0{,}682617}<0{,}001$,
cioè $0{,}6816<x_{0}<0{,}6837$; la radice dell'equazione $e^{x}+x=0$ soddisfa
$\abs{x_{0}+0{,}567383}<0{,}001$, cioè $-0{,}5684<x_{0}<-0{,}5663$. Il numero $\pm0{,}001$ è
una \emph{stima} dell'errore, non l'errore effettivo: quest'ultimo, come mostrano le
ultime righe dell'output, è di circa $2{,}9\cdot10^{-4}$ nel primo caso e di circa
$2{,}4\cdot10^{-4}$ nel secondo. Si osservi anche
che la successione $c_{n}$ non è monotòna: si avvicina a $x_{0}$ oscillando, ma la
\eqref{eq:bisezione-errore} ne controlla l'errore a ogni passo.

\begin{osservazione}[Quanti passi servono?]\label{oss:passi-bisezione}
Per ottenere $\abs{c_{n}-x_{0}}<\tau$ basta, per la \eqref{eq:bisezione-errore}, che
$\frac{b-a}{2^{n+1}}<\tau$, cioè $2^{n+1}>\frac{b-a}\tau$. Poiché $2^{10}=1024\approx10^{3}$,
ogni dieci dimezzamenti si guadagnano circa tre cifre decimali esatte: il metodo è
lento ma \emph{garantito}, e il numero di passi necessari si conosce in anticipo, prima
di cominciare il calcolo. È questa robustezza a renderlo ancora oggi il metodo di
riferimento, spesso combinato con metodi più veloci che da soli potrebbero non
convergere.
\end{osservazione}

\begin{python}[Gli zeri con la libreria \texttt{scipy}]
import math
from scipy.optimize import bisect, brentq

f = lambda x: math.exp(x) + x

# entrambe vogliono [a, b] con f(a), f(b) di segno discorde
x_b, info_b = bisect(f, -1, 0, xtol=1e-14, full_output=True)
x_r, info_r = brentq(f, -1, 0, xtol=1e-14, full_output=True)
print(f"bisect: {x_b:.14f}  ({info_b.iterations} iterazioni)")
print(f"brentq: {x_r:.14f}  ({info_r.iterations} iterazioni)")
try:
    bisect(f, 0, 1)            # f(0) = 1, f(1) = e + 1: stesso segno
except ValueError as err:
    print("errore:", err)
\end{python}
\begin{risultato}
bisect: -0.56714329040978  (47 iterazioni)
brentq: -0.56714329040978  (6 iterazioni)
errore: f(a) and f(b) must have different signs
\end{risultato}

La funzione \texttt{bisect} è esattamente il metodo della dimostrazione: con
$b-a=1$ e tolleranza $10^{-14}$ servono $47$ dimezzamenti, perché $2^{47}\approx1{,}4\cdot10^{14}$.
La funzione \texttt{brentq} (metodo di Brent) combina la bisezione con interpolazioni
più veloci e arriva allo stesso risultato in pochissime iterazioni, conservando la
garanzia di convergenza. Entrambe, come il teorema, \emph{pretendono} che $f$ cambi
segno agli estremi: senza questa ipotesi non garantiscono nulla e si rifiutano di
partire.

\subsection{Il ruolo dell'assioma di completezza}

Chiudiamo il paragrafo con un'osservazione sull'assioma di completezza
(\cref{ass:completezza}). Abbiamo utilizzato tale assioma nella dimostrazione del teorema
dell'esistenza degli zeri, in particolare nell'affermazione che la successione $a_{n}$,
essendo monotòna e limitata, risulta convergente. Ciò è essenziale; infatti,
nell'ambito dei numeri razionali $\Q$, dove non è verificato l'assioma di completezza,
non vale nemmeno il teorema dell'esistenza degli zeri. Ad esempio, la funzione
$f(x)=x^{2}-2$, considerata sull'insieme dei razionali $\{x\in\Q : 0\le x\le2\}$, è
continua e verifica $f(0)<0$, $f(2)>0$, ma non si annulla mai, perché $\sqrt2$ non è
razionale (\cref{teo:radice2}). Le successioni $a_{n}$, $b_{n}$ della bisezione sarebbero
ancora ben definite e razionali, si stringerebbero attorno al «buco» lasciato da
$\sqrt2$, ma non avrebbero un limite in $\Q$. Allo stesso modo, nella seconda
dimostrazione, l'insieme $E=\{x\in\Q : 0\le x\le2,\ x^{2}<2\}$ è limitato, ma non ha
estremo superiore in $\Q$ (\cref{teo:Q-non-completo}). L'assioma di completezza è essenziale anche
negli altri teoremi di esistenza: nel teorema di Weierstrass, o, come già detto, nel
teorema sull'esistenza del limite per le successioni monotòne.
